\subsection{Incidencia cúbica y realización lorentziana}

La completación cúbica aporta una relación entre interior y frontera.
Su geometría puede desarrollarse mediante la incidencia de las seis caras
con los ocho vértices. Etiquetemos las caras por \((i,\varepsilon)\),
\(i\in\{1,2,3\}\), \(\varepsilon\in\{+1,-1\}\), y los vértices por
\(\nu\in\{+1,-1\}^3\). La matriz de incidencia es
\[
 M_{(i,\varepsilon),\nu}=\mathbf1_{\{\nu_i=\varepsilon\}}.
\]
Sean \(u=(1,\ldots,1)^{\mathsf T}\in\RR^6\), \(R_F\) la oposición de
caras y \(P_u=uu^{\mathsf T}/6\), \(P_-=(I-R_F)/2\).

\begin{proposition}[Cociente de incidencia de dimensión cuatro]
El Gramiano de incidencia satisface
\[
 MM^{\mathsf T}=12P_u+4P_-.
\]
Sus autovalores \(12,4,0\) tienen multiplicidades \(1,3,2\).
En consecuencia,
\[
 Q_\square=\RR^6/\ker(MM^{\mathsf T})
 \simeq\RR u\oplus E_-,\qquad E_-=\operatorname{im}P_-.
\]
\end{proposition}
\begin{proof}
Una cara contiene cuatro vértices; dos caras opuestas tienen intersección
vacía, y dos caras distintas no opuestas comparten dos vértices.
La matriz \(12P_u+4P_-\) tiene precisamente esas entradas.
El espacio uniforme tiene dimensión uno; las diferencias entre caras
opuestas generan un espacio de dimensión tres, ortogonal al uniforme.
Su complemento tiene dimensión dos y pertenece al núcleo.
\end{proof}

La descomposición \(1+3\) procede de la incidencia. La elección temporal
se declara mediante la forma bilineal
\[
 B_\square([x],[y])=\langle x,(P_--P_u)y\rangle.
\]
Está bien definida en el cociente, porque ambos proyectores anulan el
núcleo. En la base
\[
 t=[u]/\sqrt6,\qquad
 d_i=[e_{i,+}-e_{i,-}]/\sqrt2,
\]
su matriz es \(\operatorname{diag}(-1,1,1,1)\). Quedan explicitados
el dominio, la descomposición y la convención de signo que producen
la realización de Minkowski. El Gramiano positivo determina el cociente;
la asignación temporal determina la forma indefinida sobre él.

El operador
\[
 W_\square=\sqrt6P_u+\sqrt2P_-,\qquad
 W_\square e_{i,\varepsilon}=t+\varepsilon d_i
\]
envía las seis etiquetas de cara a direcciones nulas, puesto que
\(B_\square(t+\varepsilon d_i,t+\varepsilon d_i)=-1+1=0\).
La inversión de una dirección y la oposición de caras disponen así
de una realización geométrica concreta.

\subsubsection{Soldadura, inversión y subdivisión}

La realización celular utiliza un transporte invertible \(U_m(a)\),
que conserva la forma lorentziana, entre las fibras de los extremos de
una arista dual orientada \(a\), una etiqueta
de cara \(\lambda_m(a)\) y una orientación \(\sigma_m(a)\).
Con \(\ell_m=\ell_0/9^m\), la aplicación de soldadura es
\[
 \beta_m(a)=\ell_m\sigma_m(a)
 W_{\square,s(a)}e_{\lambda_m(a)}.
\]
La inversión compatible satisface
\[
 \beta_m(\bar a)=-U_m(a)\beta_m(a).
\]
Identifiquemos las fibras del origen en dos niveles consecutivos mediante
la isometría de refinamiento declarada. Si los nueve subsegmentos
\(a_1,\ldots,a_9\), de nivel \(m+1\), conservan la etiqueta
y la orientación al ser transportados a esa fibra común, entonces
\begin{equation}
 \beta_m(a)=\sum_{j=1}^9
 U_{m+1}(a_1\cdots a_{j-1})^{-1}\beta_{m+1}(a_j).
 \label{eq:soldadura-refinamiento}
\end{equation}
Cada término transportado equivale a \(\ell_m/9\) por la misma
dirección nula, lo que demuestra la identidad. La hipótesis de
compatibilidad determina qué subdivisiones realizan este refinamiento.
La ecuación reúne dirección, escala y transporte, en lugar de limitar
la correspondencia geométrica a un recuento de dimensiones.

Fijadas además la orientación y la forma lorentziana, la dualidad de
Hodge en \(\Lambda^2Q_\square^*\) satisface \(\star^2=-I\).
En dimensión cuatro, el signo es \((-1)^{2(4-2)+1}=-1\).
Se obtiene una estructura compleja real sobre las dos-formas y,
tras complejificación, los subespacios propios de valores \(+i\) y
\(-i\). La realización electromagnética utilizará este dominio
después de construir la sección electrónica y declarar su acoplamiento.

\subsection{Compatibilidad entre las reducciones modular y posicional}

La base decimal por bloques interviene mediante una división euclídea
que conserva simultáneamente residuo y cociente:
\[
 n=1000q+r,\qquad 0\le r<1000.
\]
Puesto que \(1000\equiv1\pmod9\) y \(1000\equiv1\pmod3\),
\[
 n\equiv q+r\pmod9,\qquad n\equiv q+r\pmod3.
\]
Por tanto, la reducción de un bloque requiere la contribución del
cociente para recuperar la clase del entero. Los números \(0\) y
\(1000\) tienen el mismo resto por \(1000\), pero clases diferentes
por \(9\) y por \(3\). Ésta es una razón operatoria para conservar
la memoria durante el cambio de resolución.

El refinamiento ternario y la determinación de un prefijo decimal
pueden coordinarse mediante un único residuo entero. Si \(N/3^L\)
es un extremo de cilindro y \(D/1000^r\) un extremo decimal, definimos
\[
 A=1000^rN-D3^L.
\]
Al añadir seis símbolos ternarios, \(N'=729N+\nu\),
\(L'=L+6\), con \(0\le\nu<729\). Si el prefijo decimal permanece
fijo, se obtiene
\[
 A'=729A+\nu1000^r.
\]
Si se incorpora además un bloque decimal \(\delta\), entonces
\(D'=1000D+\delta\), \(r'=r+1\), y
\begin{equation}
 A'=729000A+\nu1000^{r+1}-729\delta3^L.
 \label{eq:residuo-dos-bases}
\end{equation}
Ambas identidades se demuestran sustituyendo las definiciones. La
comparación de extremos se reduce así a operaciones enteras que
retienen la contribución de cada prolongación. Esta compatibilidad
es la que permite refinar una región antes de evaluar su coordenada real.
